\documentclass[11pt]{article}
\usepackage[utf8]{inputenc}
\usepackage{amsmath,amssymb,amsfonts,amsthm}
\usepackage{geometry}
\usepackage{booktabs}
\usepackage{array}

\geometry{a4paper, margin=1in}

\title{CS282 Machine Learning: Homework 1 Solutions}
\author{Jianxuan Yao}
\date{}

\begin{document}

\maketitle

\section*{Problem 1}

\subsection*{(i)}

\begin{enumerate}
    \item \textbf{Task $T$:} Predicting whether a campus shuttle trip will arrive more than five minutes late ($Y=1$) or not ($Y=0$) before departure, given the observed traffic state $X \in \{q, b\}$.
    \item \textbf{Experience $E$:} Historical trip records consisting of observed traffic states and corresponding arrival delay outcomes.
    \item \textbf{Performance measure $P$:} The classification accuracy.
    \item \textbf{Learning type:} Supervised learning.
    \item \textbf{Task category:} Binary classification.
\end{enumerate}

\subsection*{(ii)}

$$
\mathbb{P}(X=q) = \mathbb{P}(X=q, Y=0) + \mathbb{P}(X=q, Y=1) = 0.45 + 0.05 = 0.50
$$

$$
\mathbb{P}(X=b) = \mathbb{P}(X=b, Y=0) + \mathbb{P}(X=b, Y=1) = 0.15 + 0.35 = 0.50
$$

$$
\mathbb{P}(Y=1 \mid X=q) = \frac{\mathbb{P}(X=q, Y=1)}{\mathbb{P}(X=q)} = \frac{0.05}{0.50} = 0.10
$$

$$
\mathbb{P}(Y=1 \mid X=b) = \frac{\mathbb{P}(X=b, Y=1)}{\mathbb{P}(X=b)} = \frac{0.35}{0.50} = 0.70
$$

$$
E_{\text{out}}(h_0) = \mathbb{P}(Y=1) = \mathbb{P}(X=q, Y=1) + \mathbb{P}(X=b, Y=1) = 0.05 + 0.35 = 0.40
$$

$$
E_{\text{out}}(h_1) = \mathbb{P}(X=q, Y=1) + \mathbb{P}(X=b, Y=0) = 0.05 + 0.15 = 0.20
$$

$h_1$ minimizes the population risk within $\mathcal{H}$.

\subsection*{(iii)}

\begin{enumerate}
    \item $(q, 0)$: $h_0(q)=0$ (correct), $h_1(q)=0$ (correct)
    \item $(q, 0)$: $h_0(q)=0$ (correct), $h_1(q)=0$ (correct)
    \item $(b, 0)$: $h_0(b)=0$ (correct), $h_1(b)=1$ (incorrect)
    \item $(b, 0)$: $h_0(b)=0$ (correct), $h_1(b)=1$ (incorrect)
    \item $(b, 1)$: $h_0(b)=0$ (incorrect), $h_1(b)=1$ (correct)
\end{enumerate}

The empirical risks are:

$$
E_{\text{in}}(h_0) = \frac{1}{5} \sum_{i=1}^5 \mathbf{1}\{h_0(x_i) \neq y_i\} = \frac{1}{5} = 0.20
$$

$$
E_{\text{in}}(h_1) = \frac{1}{5} \sum_{i=1}^5 \mathbf{1}\{h_1(x_i) \neq y_i\} = \frac{2}{5} = 0.40.
$$

Because $E_{\text{in}}(h_0) < E_{\text{in}}(h_1)$, the ERM choice is $h_0$.

\subsection*{(iv)}

This will lead to data leakage. The stated task requires predicting trip delays before departure. However, the delay reason entered by the driver is generated after arrival. Therefore, the model cannot be deployed or used for its intended purpose.

\newpage

\section*{Problem 2}

\subsection*{(i)}

Define the augmented input vector $\tilde{x}_i \in \mathbb{R}^{d+1}$ and the augmented parameter vector $\tilde{w} \in \mathbb{R}^{d+1}$ as:

$$
\tilde{x}_i = \begin{bmatrix} x_i \\ 1 \end{bmatrix}, \quad \tilde{w} = \begin{bmatrix} w \\ b \end{bmatrix}
$$

\begin{enumerate}
    \item \textbf{Score:}
    
    $$
    \tilde{s}(\tilde{x}) = \tilde{w}^\top \tilde{x} = \begin{bmatrix} w^\top & b \end{bmatrix} \begin{bmatrix} x \\ 1 \end{bmatrix} = w^\top x + b = s(x)
    $$

    \item \textbf{Update Condition:}
    
    $$
    y_i \tilde{w}^\top \tilde{x}_i = y_i (w^\top x_i + b) \le 0
    $$

    Thus, an update is triggered in the augmented space if and only if it is triggered in the separated model.
    
    \item \textbf{Initialization:}
    
    The initial conditions $w^{(0)} = \mathbf{0}_d$ and $b^{(0)} = 0$ correspond directly to $\tilde{w}^{(0)} = \begin{bmatrix} \mathbf{0}_d \\ 0 \end{bmatrix} = \mathbf{0}_{d+1}$.
    
    \item \textbf{Parameter Update:}
    
    When triggered, the augmented parameter update is:
    
    $$
    \tilde{w} \leftarrow \tilde{w} + y_i \tilde{x}_i = \begin{bmatrix} w \\ b \end{bmatrix} + y_i \begin{bmatrix} x_i \\ 1 \end{bmatrix} = \begin{bmatrix} w + y_i x_i \\ b + y_i \end{bmatrix},
    $$

    which exactly matches the separate updates $w \leftarrow w + y_i x_i$ and $b \leftarrow b + y_i$.
\end{enumerate}

Therefore, the two formulations are identical under the same visit sequence.

\subsection*{(ii)}

\begin{center}
\begin{tabular}{cccccc}
\toprule
Visit $i$ & Example $(x_i, y_i)$ & Margin before update $y_i(wx_i + b)$ & Parameter Update & Resulting $(w, b)$ \\
\midrule
1 & $(1, -1)$ & $-1 \cdot (0 \cdot 1 + 0) = 0 \le 0$  & $w \leftarrow -1$ & $(-1, -1)$ \\
  &          &                                     & $b \leftarrow -1$     &   &        \\
2 & $(2, -1)$ & $-1 \cdot (-1 \cdot 2 - 1) = +3 > 0$  & None                             & $(-1, -1)$ \\
3 & $(3, +1)$ & $+1 \cdot (-1 \cdot 3 - 1) = -4 \le 0$ & $w \leftarrow 2$  & $(2, 0)$   \\
  &          &                                      & $b \leftarrow 0$ &         &           \\
4 & $(4, +1)$ & $+1 \cdot (2 \cdot 4 + 0) = +8 > 0$    & None                             & $(2, 0)$   \\
\bottomrule
\end{tabular}
\end{center}

The final parameters are $(w, b) = (2, 0)$. Checking all four examples:

\begin{enumerate}
    \item $i=1$: $y_1(2 \cdot 1 + 0) = -1 \cdot 2 = -2 \le 0$ (incorrect)
    \item $i=2$: $y_2(2 \cdot 2 + 0) = -1 \cdot 4 = -4 \le 0$ (incorrect)
    \item $i=3$: $y_3(2 \cdot 3 + 0) = +1 \cdot 6 = +6 > 0$ (correct)
    \item $i=4$: $y_4(2 \cdot 4 + 0) = +1 \cdot 8 = +8 > 0$ (correct)
\end{enumerate}

\subsection*{(iii)}

Choose a decision boundary at $x = 2.5$, given by $w = 2, b = -5$:

$$
i=1: \quad y_1(w x_1 + b) = -1 \cdot (2(1) - 5) = +3 > 0
$$

$$
i=2: \quad y_2(w x_2 + b) = -1 \cdot (2(2) - 5) = +1 > 0
$$

$$
i=3: \quad y_3(w x_3 + b) = +1 \cdot (2(3) - 5) = +1 > 0
$$

$$
i=4: \quad y_4(w x_4 + b) = +1 \cdot (2(4) - 5) = +3 > 0
$$

If $b=0$, the margin is $y_i (w x_i)$.

For $x_1 = 1, y_1 = -1$, strict separation requires $-1 \cdot (w \cdot 1) > 0 \implies w < 0$; But for $x_3 = 3, y_3 = +1$, strict separation requires $+1 \cdot (w \cdot 3) > 0 \implies w > 0$.

Since no real number satisfies both $w < 0$ and $w > 0$, no such $w$ exists.

Fixing $b=0$ forces the decision boundary hyperplanes to pass through the origin. With a free bias $b$, the decision boundary can be translated anywhere.

\subsection*{(iv)}
Any threshold $t \in [2, 3)$ correctly classifies all four training samples.

\begin{enumerate}
    \item Choose $t_1 = 2.2$: here $2.5 > 2.2 \implies f_{2.2}(2.5) = +1$
    \item Choose $t_2 = 2.7$: here $2.5 \le 2.7 \implies f_{2.7}(2.5) = -1$
\end{enumerate}

The training data don't contain $(2, 3)$, so in this invertal we lack information to determine.

\newpage

\section*{Problem 3}

\subsection*{(i)}

$$
J(\theta) = \frac{1}{2n} \|X\theta - y\|_2^2 = \frac{1}{2n} \left( \theta^\top X^\top X \theta - 2 y^\top X \theta + y^\top y \right)
$$

$$
\nabla_\theta J(\theta) = \frac{1}{n} X^\top (X\theta - y)
$$

When $\nabla_\theta J(\theta) = 0$, $X^\top X \theta = X^\top y$.
\begin{enumerate}
    \item $X^\top X$ is invertible if and only if $X$ has full column rank, which means that the columns of $X$ are linearly independent.
    \item Dimensions:
    \begin{enumerate}
        \item $X^\top X \in \mathbb{R}^{2 \times 2}$
        \item $X^\top y \in \mathbb{R}^{2 \times 1}$
        \item $\nabla J \in \mathbb{R}^{2 \times 1}$
    \end{enumerate}
\end{enumerate}

\subsection*{(ii)}

$$
X = \begin{bmatrix} 1 & -2 \\ 1 & -1 \\ 1 & 1 \\ 1 & 2 \end{bmatrix}, \quad y = \begin{bmatrix} 0 \\ 1 \\ 2 \\ 5 \end{bmatrix}
$$

$$
X^\top X = \begin{bmatrix} 4 & 0 \\ 0 & 10 \end{bmatrix}
$$

$$
X^\top y = \begin{bmatrix} 8 \\ 11 \end{bmatrix}
$$

$$
\hat{\theta} = (X^\top X)^{-1} X^\top y = \begin{bmatrix} 1/4 & 0 \\ 0 & 1/10 \end{bmatrix} \begin{bmatrix} 8 \\ 11 \end{bmatrix} = \begin{bmatrix} 2 \\ 1.1 \end{bmatrix}
$$

$$
r = X\hat{\theta} - y = \begin{bmatrix} -0.2 \\ -0.1 \\ 1.1 \\ -0.8 \end{bmatrix}
$$

$$
X^\top r = \begin{bmatrix} 1 & 1 & 1 & 1 \\ -2 & -1 & 1 & 2 \end{bmatrix} \begin{bmatrix} -0.2 \\ -0.1 \\ 1.1 \\ -0.8 \end{bmatrix} = \mathbf{0}
$$

$$
\Vert r \Vert_2^2 = 1.90
$$

$$
\text{MSE} = \frac{1}{n} \|r\|_2^2 = 0.475
$$

$$    
J(\hat{\theta}) = \frac{1}{2n} \|r\|_2^2 = 0.2375
$$

\subsection*{(iii)}

Having $n \ge p$ is merely a necessary condition for $Z^\top Z$ to be invertible, not a sufficient one. $rank(Z^\top Z) = rank(Z) = 2 < 3$, so $Z^\top Z$ is not invertible.

The linear model predicts:

$$
\hat{y} = a + bx + cz = a + (b + 2c)x
$$

To minimize the sum of squared errors, the effective slope and intercept must match the unique optimal values obtained in part (ii):

$$
a = 2 \quad \text{and} \quad b + 2c = 1.1
$$

Thus, the set of all minimizing parameter triples is:

$$
\mathcal{S} = \left\{ (a, b, c) \in \mathbb{R}^3 \;\middle|\; a = 2, \; b + 2c = 1.1 \right\}
$$

\begin{enumerate}
    \item Setting $c = 0 \implies (a, b, c) = (2, 1.1, 0)$.
    \item Setting $c = 0.55 \implies (a, b, c) = (2, 0, 0.55)$.
\end{enumerate}

For a future inputs satisfying $z = 2x$, the predictions are identical. The predicted value is:

$$
\hat{y} = a + bx + cz = a + (b + 2c)x = 2 + 1.1x.
$$

\newpage

\section*{Problem 4}

\subsection*{(i)}

$$
L(\theta) = \prod_{i=1}^n p_i^{t_i} (1 - p_i)^{1 - t_i}
$$

$$
J(\theta)
= -\frac{1}{n} \ln L(\theta)
= -\frac{1}{n} \sum_{i=1}^n \left[ t_i \ln p_i + (1 - t_i) \ln(1 - p_i) \right]
$$

\begin{enumerate}
    \item If $t_i = 1$, then $y_i = 2(1) - 1 = +1$:
    
    $$
    -\ln p_i = -\ln \left( \frac{1}{1 + e^{-\theta^\top x_i}} \right) = \ln(1 + e^{-\theta^\top x_i}) = \ln(1 + e^{-y_i \theta^\top x_i})
    $$
    
    \item If $t_i = 0$, then $y_i = 2(0) - 1 = -1$:
    
    $$
    -\ln(1 - p_i) = -\ln \left( \frac{1}{1 + e^{\theta^\top x_i}} \right) = \ln(1 + e^{\theta^\top x_i}) = \ln(1 + e^{-y_i \theta^\top x_i})
    $$
\end{enumerate}

Combining both cases into a single sum:

$$
J(\theta) = \frac{1}{n} \sum_{i=1}^n \ln(1 + e^{-y_i \theta^\top x_i})
$$

\subsection*{(ii)}

$$
\nabla_\theta p_i = \nabla_\theta \sigma(\theta^\top x_i) = \sigma'(\theta^\top x_i) x_i = p_i(1 - p_i) x_i
$$

$$
\begin{aligned}
&\nabla_\theta \left( -[t_i \ln p_i + (1 - t_i) \ln (1 - p_i)] \right) \\
=& -\left( \frac{t_i}{p_i} - \frac{1 - t_i}{1 - p_i} \right) \nabla_\theta p_i \\
=& -\left( \frac{t_i(1 - p_i) - (1 - t_i)p_i}{p_i(1 - p_i)} \right) p_i(1 - p_i) x_i \\
=& (p_i - t_i) x_i
\end{aligned}
$$

$$
\nabla J(\theta) = \frac{1}{n} \sum_{i=1}^n (p_i - t_i) x_i = \frac{1}{n} X^\top (p - t)
$$

$$
d (\nabla J(\theta))
= \frac{1}{n} X^\top d(\sigma(X \theta))
= \frac{1}{n} X^\top \sigma'(X \theta) \odot X d \theta
= \frac{1}{n} X^\top \operatorname{diag}(\sigma'(X \boldsymbol{\theta})) X d\boldsymbol{\theta}
$$

$$
\nabla^2 J(\boldsymbol{\theta}) = \frac{\partial \nabla J(\boldsymbol{\theta})}{\partial \boldsymbol{\theta}^\top} = \frac{1}{n} X^\top \operatorname{diag}\left(\sigma(X\boldsymbol{\theta}) \odot (\mathbf{1} - \sigma(X\boldsymbol{\theta}))\right) X
$$

$$
D_{ii} = p_i(1 - p_i) = \sigma(\theta^\top x_i)(1 - \sigma(\theta^\top x_i))
$$

\subsection*{(iii)}

For any non-zero vector $v \in \mathbb{R}^d$, consider the quadratic form:

$$
v^\top H v = v^\top \left( \frac{1}{n} X^\top D X \right) v = \frac{1}{n} (Xv)^\top D (Xv)
$$

Let $z = Xv \in \mathbb{R}^n$. Then:

$$
v^\top H v = \frac{1}{n} \sum_{i=1}^n D_{ii} z_i^2 = \frac{1}{n} \sum_{i=1}^n p_i(1 - p_i) z_i^2
$$

For any finite $\theta$ and finite $x_i$, the sigmoid output satisfies $p_i = \sigma(\theta^\top x_i) \in (0, 1)$, which implies $D_{ii} = p_i(1 - p_i) > 0$.

Since each term in the sum is the product of a strictly positive quantity $D_{ii}$ and a non-negative quantity $z_i^2 \ge 0$, we have:

$$
v^\top H v \ge 0 
$$

Therefore, $H$ is semidefinite.

If the columns of $X$ are linearly independent, the null space of $X$ contains only the zero vector:

$$
Xv = 0 \iff v = 0.
$$

Hence, for any $v \neq 0$, the vector $z = Xv \neq 0$, which means at least one component satisfies $z_k^2 > 0$. Since all diagonal elements satisfy $D_{ii} > 0$, we have

$$
v^\top H v = \frac{1}{n} \sum_{i=1}^n D_{ii} z_i^2 > 0.
$$

Therefore, $H$ is strictly positive definite.

\subsection*{(iv)}

$$
X = \begin{bmatrix} 1 & -2 \\ 1 & 0 \\ 1 & 2 \end{bmatrix}
,\quad
t = \begin{bmatrix} 0 \\ 1 \\ 1 \end{bmatrix}
$$

At $\theta = (0, 0)^\top$, we have $\theta^\top x_i = 0$ for all $i$, so

$$
p = \begin{bmatrix} 1/2 \\ 1/2 \\ 1/2 \end{bmatrix}
, \quad
p - t = \begin{bmatrix} 1/2 \\ -1/2 \\ -1/2 \end{bmatrix}
$$

$$
\nabla J(\mathbf{0})
= \frac{1}{3} X^\top (p - t)
= \frac{1}{3} \begin{bmatrix} 1 & 1 & 1 \\ -2 & 0 & 2 \end{bmatrix} \begin{bmatrix} 1/2 \\ -1/2 \\ -1/2 \end{bmatrix}
= \begin{bmatrix} -1/6 \\ -2/3 \end{bmatrix}.
$$

For the Hessian, $D_{ii} = \frac{1}{2} \left( 1 - \frac{1}{2} \right) = \frac{1}{4}$ for all $i$, so $D = \frac{1}{4} I_3$:

$$
X^\top X = \begin{bmatrix} 3 & 0 \\ 0 & 8 \end{bmatrix} \implies H = \frac{1}{3} X^\top \left(\frac{1}{4} I_3\right) X = \frac{1}{12} \begin{bmatrix} 3 & 0 \\ 0 & 8 \end{bmatrix} = \begin{bmatrix} 1/4 & 0 \\ 0 & 2/3 \end{bmatrix}.
$$

\newpage

\section*{Problem 5}

\subsection*{(i)}

Let $\hat{Y} \in \{0, 1\}$ be the predicted decision.

$$
\begin{aligned}
\mathbb{E}[\text{Cost} \mid X=x, \hat{Y}=0] &= p(x) C_{FN} + (1 - p(x)) \cdot 0 = p(x) C_{FN}, \\
\mathbb{E}[\text{Cost} \mid X=x, \hat{Y}=1] &= p(x) \cdot 0 + (1 - p(x)) C_{FP} = (1 - p(x)) C_{FP}.
\end{aligned}
$$

To minimize the conditional expected cost, predict $\hat{Y}=1$ if and only if:

$$
\mathbb{E}[\text{Cost} \mid \hat{Y}=1] \le \mathbb{E}[\text{Cost} \mid \hat{Y}=0] \iff (1 - p(x)) C_{FP} \le p(x) C_{FN},
$$

$$
p(x) \ge \frac{C_{FP}}{C_{FP} + C_{FN}}.
$$

$$
\tau^* = \frac{1}{1 + 4} = 0.2.
$$

Predict $1$ if $p(x) \ge 0.2$, and $0$ otherwise.

\subsection*{(ii)}

\textbf{Case 1: Threshold $\tau = 0.5$}

\begin{enumerate}
    \item Predictions:
    
    $$
    \hat{y} = (1, 1, 1, 0, 0, 0)
    $$

    \item Counts:
    
    $$
    \text{TP} = 2
    , \quad \text{FP} = 1
    , \quad \text{FN} = 1
    , \quad \text{TN} = 2
    $$

    \item Metrics:
    
    $$
    \begin{aligned}
    \text{Accuracy} &= \frac{\text{TP} + \text{TN}}{P + N} \approx 66.67\%, \\
    \text{Precision} &= \frac{\text{TP}}{\text{TP} + \text{FP}} \approx 66.67\%, \\
    \text{Recall} &= \frac{\text{TP}}{\text{TP} + \text{FN}} \approx 66.67\%, \\
    \text{F1} &= \frac{2 \cdot \text{Precision} \cdot \text{Recall}}{\text{Precision} + \text{Recall}} \approx 0.6667.
    \end{aligned}
    $$

    \item Total Observed Cost:
    
    $$
    \text{Cost} = 5
    $$
\end{enumerate}

\textbf{Case 2: Threshold $\tau = 0.2$}

\begin{enumerate}
    \item Predictions:
    
    $$
    \hat{y} = (1, 1, 1, 1, 1, 0)
    $$

    \item Counts:
    
    $$
    \text{TP} = 3
    , \quad \text{FP} = 2
    , \quad \text{FN} = 0,
    \quad \text{TN} = 1
    $$

    \item Metrics:
    
    $$
    \begin{aligned}
    \text{Accuracy} &= \frac{\text{TP} + \text{TN}}{P + N} \approx 66.67\%, \\
    \text{Precision} &= \frac{\text{TP}}{\text{TP} + \text{FP}} \approx 60.00\%, \\
    \text{Recall} &= \frac{\text{TP}}{\text{TP} + \text{FN}} \approx 100.00\%, \\
    \text{F1} &= \frac{2 \cdot \text{Precision} \cdot \text{Recall}}{\text{Precision} + \text{Recall}} \approx 0.75.
    \end{aligned}
    $$

    \item Total Observed Cost:
    
    $$
    \text{Cost} = 2
    $$
\end{enumerate}

Both rules share an identical accuracy of $66.67\%$, yet the threshold $\tau=0.2$ cuts total cost from 5 to 2.

This demonstrates that choosing a decision threshold based solely on accuracy ignores the real-world misclassification costs.

\subsection*{(iii)}

\begin{center}
\begin{tabular}{ccccccc}
\toprule
Rank & Score & Label ($Y$) & Threshold range $\tau$ & Cumulative (FP, TP) & Coordinate (FPR, TPR) \\
\midrule
-- & --   & -- & $\tau > 0.90$          & $(0, 0)$ & $(0, 0)$ \\
1  & 0.90 & 1  & $0.70 < \tau \le 0.90$ & $(0, 1)$ & $(0, 1/3)$ \\
2  & 0.70 & 0  & $0.60 < \tau \le 0.70$ & $(1, 1)$ & $(1/3, 1/3)$ \\
3  & 0.60 & 1  & $0.40 < \tau \le 0.60$ & $(1, 2)$ & $(1/3, 2/3)$ \\
4  & 0.40 & 0  & $0.30 < \tau \le 0.40$ & $(2, 2)$ & $(2/3, 2/3)$ \\
5  & 0.30 & 1  & $0.10 < \tau \le 0.30$ & $(2, 3)$ & $(2/3, 1)$ \\
6  & 0.10 & 0  & $\tau \le 0.10$        & $(3, 3)$ & $(1, 1)$ \\
\bottomrule
\end{tabular}
\end{center}

The successive ROC coordinates are:

$$
(0, 0)
\to \left(0, \frac{1}{3}\right)
\to \left(\frac{1}{3}, \frac{1}{3}\right)
\to \left(\frac{1}{3}, \frac{2}{3}\right)
\to \left(\frac{2}{3}, \frac{2}{3}\right)
\to \left(\frac{2}{3}, 1\right)
\to (1, 1)
$$

$$
\text{AUC}
= \left(\frac{1}{3} - 0\right) \cdot \frac{1}{3} + \left(\frac{2}{3} - \frac{1}{3}\right) \cdot \frac{2}{3} + \left(1 - \frac{2}{3}\right) \cdot 1
= \frac{1}{9} + \frac{2}{9} + \frac{3}{9} = \frac{6}{9}
= \frac{2}{3}
$$

\subsection*{(iv)}

Neither the ROC curve nor the AUC changes. The mapping $g(p) = p^2$ is strictly monotonically increasing on $[0, 1]$. It would not change the relative ordering of the scores, and thus the ROC curve remains unchanged.

Computing the transformed scores $q_i = p_i^2$:

$$
\begin{aligned}
q_1 &= 0.9^2 = 0.81 \ge 0.2 \implies 1 \\
q_2 &= 0.7^2 = 0.49 \ge 0.2 \implies 1 \\
q_3 &= 0.6^2 = 0.36 \ge 0.2 \implies 1 \\
q_4 &= 0.4^2 = 0.16 < 0.2 \implies 0 \\
q_5 &= 0.3^2 = 0.09 < 0.2 \implies 0 \\
q_6 &= 0.1^2 = 0.01 < 0.2 \implies 0
\end{aligned}
$$

$$
\hat{y} = (1, 1, 1, 0, 0, 0)
$$

$$
\text{Cost} = 5.
$$

Because $p \ge 0$, we have:

$$
p \ge 0.2 \iff p^2 \ge 0.2^2 = 0.04.
$$

Therefore, the corresponding threshold on $q$ is $\tau_q = 0.04$.
\end{document}